\section{Selecting an embedding by rank support}
\label{sec:support}

The preceding section supplies a nonprojective representation of an
elementary abelian group whose restrictions to all proper subgroups are
projective.  We now choose an embedding of that group in the unipotent
subgroup of $\mathrm{SL}_2$ so that its tensor product with the restricted
Steinberg module remains nonprojective.  This choice will make the higher
Frobenius functor exact.

Let $k$ be algebraically closed of characteristic $p>0$, let
$E_r=(\mathbb Z/p\mathbb Z)^r$, and fix generators $g_1,\ldots,g_r$.
For a finite-dimensional $kE_r$-module $D$, define
\[
 \operatorname{Supp}(D)=
 \left\{[\alpha_1:\cdots:\alpha_r]\in\mathbb P^{r-1}(k):
 D\text{ is not free over }k[t]/(t^p),\quad
 t\mapsto\sum_{i=1}^r\alpha_i(g_i-1)\right\}.
\]
This is the projectivized rank variety.  The rank-variety projectivity
and tensor-product theorems state that
\begin{equation}\label{eq:support-properties}
 \operatorname{Supp}(D)=\varnothing\iff D\text{ is projective},
 \qquad
 \operatorname{Supp}(D\otimes D')
 =\operatorname{Supp}(D)\cap\operatorname{Supp}(D').
\end{equation}
These statements hold with the indicated coordinates for every prime;
see \cite[Sections~1.9 and~1.12]{Benson2017} and
\cite[Section~2.3]{CoulembierFlakeSteinberg}.

\begin{lemma}\label{lem:support-independent}
Suppose that $M$ is a nonprojective finite-dimensional $kE_r$-module and
that $M|_H$ is projective for every proper subgroup $H<E_r$.
Every point $[\alpha_1:\cdots:\alpha_r]$ of $\operatorname{Supp}(M)$ has
$\mathbb F_p$-linearly independent coordinates.
\end{lemma}

\begin{proof}
For a nonzero linear form $c:E_r\to\mathbb F_p$, write $c(g_i)=c_i$ and
put $H=\ker c$.  The tensor identity gives
\[
 M\otimes\operatorname{Ind}_H^{E_r}k
 \cong\operatorname{Ind}_H^{E_r}(M|_H),
\]
which is projective.  We compute the support of
$N=\operatorname{Ind}_H^{E_r}k$ explicitly.  This is the regular
representation of $E_r/H$, identified with $k[z]/(z^p)$, where a
generator of the quotient acts by $1+z$.  Thus
$\sum_i\alpha_i(g_i-1)$ acts by multiplication by
\[
 h_\alpha(z)=\sum_i\alpha_i\big((1+z)^{c_i}-1\big),
\]
with $c_i$ represented by an integer in $\{0,\ldots,p-1\}$.
Its linear coefficient is $\sum_i c_i\alpha_i$.
If this coefficient is nonzero, substitution $t\mapsto h_\alpha(z)$
is an isomorphism of truncated polynomial algebras, so $N$ is free
over $k[t]/(t^p)$.  If it is zero, then $h_\alpha\in(z^2)$ and
$h_\alpha^{p-1}=0$, which is impossible on a free module of dimension
$p$.  Consequently
\[
 \operatorname{Supp}(N)
 =\{[\alpha]:\textstyle\sum_i c_i\alpha_i=0\}.
\]
Equation~\eqref{eq:support-properties} shows that this hyperplane misses
$\operatorname{Supp}(M)$.  Varying the nonzero vector
$(c_1,\ldots,c_r)\in\mathbb F_p^r$ proves the assertion.
\end{proof}

Every support point of the minimal-rank module therefore has
$\mathbb F_p$-linearly independent coordinates.
To make one of them survive after tensoring with a restricted Steinberg
module, we next compute the Steinberg support in terms of the embedding
parameters.

\begin{lemma}\label{lem:steinberg-support}
Let $\lambda_1,\ldots,\lambda_r\in k$ be linearly independent over
$\mathbb F_p$, and restrict $\mathrm{SL}_2$-representations along
\[
 \iota_\lambda:E_r\longrightarrow\mathrm{SL}_2(k),
 \qquad
 g_i\longmapsto
 \begin{pmatrix}1&\lambda_i\\0&1\end{pmatrix}.
\]
Write $R_\lambda$ for this restriction.  Then
\begin{equation}\label{eq:steinberg-support}
 \operatorname{Supp}\big(R_\lambda(\operatorname{St}_{r-1})\big)
 =\left\{[\alpha]:
 \sum_i\alpha_i\lambda_i^{p^j}=0\quad(0\le j\le r-2)\right\}.
\end{equation}
\end{lemma}

\begin{proof}
This is the Steinberg case of
\cite[Proposition~4.2.5 and Corollary~4.2.7]{CoulembierFlakeSteinberg}.
Here is a direct calculation in the chosen coordinates.
On $\operatorname{St}_1=\operatorname{Sym}^{p-1}(k^2)$, the upper
unipotent element with parameter $a$ acts by
$\exp(aN)=\sum_{b=0}^{p-1}a^bN^b/b!$, where $N$ has one Jordan
block of length $p$.  Therefore $\sum_i\alpha_i(g_i-1)$ acts by a
polynomial in $N$ whose linear coefficient is
$\sum_i\alpha_i\lambda_i$.  As in the proof of
Lemma~\ref{lem:support-independent}, the resulting module of dimension
$p$ is free exactly when this coefficient is nonzero.
On the $j$-th Frobenius twist, the same calculation replaces
$\lambda_i$ by $\lambda_i^{p^j}$.
Now use
\[
 \operatorname{St}_{r-1}
 \cong\bigotimes_{j=0}^{r-2}\operatorname{St}_1^{(j)}
\]
and the tensor-product formula in \eqref{eq:support-properties}.
When $r=1$, $\operatorname{St}_0=k$ has support
$\mathbb P^0(k)$, agreeing with the empty system in
\eqref{eq:steinberg-support}.
\end{proof}

Thus a point $[\alpha]$ lies in the Steinberg support precisely when
the embedding parameters satisfy the equations in
\eqref{eq:steinberg-support}. The next lemma solves these equations
while keeping the parameters linearly independent over $\mathbb F_p$,
as required for an embedding.

\begin{lemma}\label{lem:moore-selection}
Let $\alpha_1,\ldots,\alpha_r\in k$ be linearly independent over
$\mathbb F_p$.  There are $\mathbb F_p$-linearly independent
$\lambda_1,\ldots,\lambda_r\in k$ such that
\begin{equation}\label{eq:selected-embedding}
 \sum_{i=1}^r\alpha_i\lambda_i^{p^j}=0
 \qquad(0\le j\le r-2).
\end{equation}
For $r=1$ the list of equations is empty.
\end{lemma}

\begin{proof}
When $r=1$, take any $\lambda_1\ne0$.  Suppose $r\ge2$.
Because $k$ is perfect, we may put
$\beta_i=\alpha_i^{p^{-(r-2)}}$.  Form the $(r-1)\times r$ matrix
\[
 B=(\beta_i^{p^s})_{\substack{0\le s\le r-2\\1\le i\le r}}.
\]
Its rows are linearly independent.  Indeed, a nontrivial row relation
would give a nonzero additive polynomial
$f(X)=\sum_{s=0}^{r-2}a_sX^{p^s}$ vanishing on the
$\mathbb F_p$-span of the $\beta_i$.  That span has $p^r$ elements,
whereas $\deg f\le p^{r-2}$.

The row space of $B$ contains no nonzero vector in $\mathbb F_p^r$.
To see this, suppose $c_i=f(\beta_i)\in\mathbb F_p$ for such a
polynomial $f$.  The additive polynomial $f^p-f$ then vanishes on
the same $p^r$ elements and has degree at most $p^{r-1}$.
Thus $f^p-f=0$, forcing $f=0$ because $f$ has zero constant term.

Choose a nonzero vector $\lambda\in\ker B$.  This kernel is a line.
If $c\in\mathbb F_p^r$ satisfied $\sum_i c_i\lambda_i=0$, it would
belong to the annihilator of $\ker B$, which is the row space of $B$.
Hence $c=0$, proving the required independence of the $\lambda_i$.
The rows of $B$, in reverse order, are
$(\alpha_i^{p^{-j}})_i$ for $0\le j\le r-2$.
Raising their nullvector equations to the $p^j$-th power gives
\eqref{eq:selected-embedding}.
\end{proof}

\begin{proposition}\label{prop:support-selection}
Let $M$ be a nonprojective finite-dimensional $kE_r$-module whose
restriction to every proper subgroup is projective.  There is an
embedding $\iota_\lambda:E_r\hookrightarrow\mathrm{SL}_2(k)$ into
the upper unipotent subgroup such that
\[
 R_\lambda(\operatorname{St}_{r-1})\otimes M
 \quad\text{is nonprojective}.
\]
\end{proposition}

\begin{proof}
By \eqref{eq:support-properties}, choose
$[\alpha]\in\operatorname{Supp}(M)$.
Lemma~\ref{lem:support-independent} makes its coordinates
$\mathbb F_p$-linearly independent.
Choose $\lambda$ by Lemma~\ref{lem:moore-selection}.
Its independent coordinates make $\iota_\lambda$ injective, and
Lemma~\ref{lem:steinberg-support} places $[\alpha]$ in the support of
$R_\lambda(\operatorname{St}_{r-1})$.
The tensor-product theorem now gives the conclusion.
\end{proof}

\begin{theorem}[An exact higher Frobenius functor]
\label{thm:exact-higher-frobenius}
Let $\mathcal C$ be a finite symmetric tensor category over an
algebraically closed field $k$ of characteristic $p>0$.
There are an integer $r\ge1$ and an embedding
$E_r\hookrightarrow U(k)$ for which the construction of
Section~\ref{sec:higher} gives an exact additive strong symmetric
monoidal functor
\[
 \Phi:\mathcal C\longrightarrow
 \mathcal C\boxtimes\operatorname{Ver}_{p^r}(k).
\]
It is given by $\Phi(X)=\Theta_{\mathcal C}(X^{\otimes p^r})$, where
$\Theta$ is the functor in \eqref{eq:theta} for the chosen embedding.
The functor $\Phi$ satisfies
$\Phi(af)=a^{p^r}\Phi(f)$, and its evaluation at every projective
$Q\in\mathcal C$ is given by the natural isomorphism
\eqref{eq:phi-evaluation}.
\end{theorem}

\begin{proof}
Choose a projective generator $P$ and the least rank $r$ of
Proposition~\ref{prop:first-nonprojective}. Then
\[
 M=\operatorname{Hom}_{\mathcal C}(P,P^{\otimes p^r})
\]
is nonprojective and has projective restriction to every proper subgroup
of $E_r$. Proposition~\ref{prop:support-selection} supplies an embedding
$\iota_\lambda$ for which
$R_\lambda(\operatorname{St}_{r-1})\otimes M$ remains nonprojective.
Theorem~\ref{thm:restricted-tilting} gives
$R_\lambda(\operatorname{St}_{r-1})\otimes M\cong R_\lambda(T)$
for a tilting module $T$. Lemma~\ref{lem:theta-detection} therefore gives
\[
 0\ne\Theta\big(R_\lambda(T)\big)
 \cong\Theta\big(R_\lambda(\operatorname{St}_{r-1})\big)
       \otimes\Theta(M).
\]
In particular $\Theta(M)\ne0$. Proposition~\ref{prop:higher-exactness},
with both projectives equal to $P$, proves that $\Phi$ is exact.
Its remaining properties and the projective-evaluation identity were
established in Section~\ref{sec:higher}. All steps include $p=2$ and $r=1$.
\end{proof}
